\documentclass[letterpaper,10pt,english]{article}
\usepackage{%
amsfonts,%
amsmath,%
amssymb,%
amsthm,%
babel,%
bbm,%
%biblatex,%
caption,%
centernot,%
color,%
enumerate,%
%enumitem,%
epsfig,%
epstopdf,%
etex,%
fancybox,%
framed,%
fullpage,%
%geometry,%
graphicx,%
hyperref,%
latexsym,%
mathptmx,%
mathtools,%
multicol,%
paralist,%
pgf,%
pgfplots,%
pgfplotstable,%
pgfpages,%
proof,%
psfrag,%
%subfigure,%
tcolorbox,%
tfrupee,%
tikz,%
times,%
%ulem,%
url,%
xcolor,%
mathpazo,%
}

\definecolor{shadecolor}{gray}{.95}%{rgb}{1,0,0}
\usepackage[margin=1in,top=0.75in]{geometry}
\usepackage[mathscr]{eucal}
\usepgflibrary{shapes}
\usepgfplotslibrary{fillbetween}
\usetikzlibrary{%
arrows,%
backgrounds,%
chains,%
decorations.pathmorphing,% /pgf/decoration/random steps | erste Graphik
decorations.text,% 
matrix,%
positioning,% wg. " of "
fit,%
patterns,%
petri,%
plotmarks,%
scopes,%
shadows,%
shapes.misc,% wg. rounded rectangle
shapes.arrows,%
shapes.callouts,%
shapes%
}

%\pgfplotsset{compat=newest} %<------ Here
\pgfplotsset{compat=1.11} %<------ Or use this one

\theoremstyle{plain}
\newtheorem{thm}{Theorem}[section]
\newtheorem{lem}[thm]{Lemma}
\newtheorem{prop}[thm]{Proposition}
\newtheorem{cor}[thm]{Corollary}
\newtheorem{clm}[thm]{Claim}

\theoremstyle{definition}
\newtheorem{defn}[thm]{Definition}
\newtheorem{conj}[thm]{Conjecture}
\newtheorem{exmp}[thm]{Example}
\newtheorem{axiom}[thm]{Axiom}
\newtheorem{assum}[thm]{Assumption}
\newtheorem{exerc}[thm]{Exercise}
\newtheorem*{defin}{Definition}

\theoremstyle{remark}
\newtheorem{rem}{Remark}
\newtheorem{note}{Note}

\newcommand{\iid}{\emph{i.i.d.}\ }
\newcommand{\Cov}{\operatorname{Cov}}
%\newcommand{\det}{\operatorname{det}}
%\newcommand{\Real}{\mathbb{R}}
\newcommand{\tr}{\operatorname{tr}}
\newcommand{\Var}{\operatorname{Var}}
\newcommand{\cov}{\operatorname{cov}}

%\renewcommand{\proof}[1]{\begin{proof}#1\end{proof}}
\newcommand{\EQ}[1]{\begin{equation*}#1\end{equation*}}
\newcommand{\EQN}[1]{\begin{equation}#1\end{equation}}
\newcommand{\eq}[1]{\begin{align*}#1\end{align*}}
\newcommand{\meq}[2]{\begin{xalignat*}{#1}#2\end{xalignat*}}
\newcommand{\norm}[1]{\left\lVert#1\right\rVert}
\newcommand{\inner}[1]{\left\langle #1\right\rangle}
\newcommand{\abs}[1]{\left\lvert#1\right\rvert}
\newcommand{\expect}[1]{\mathbb{E}\left[{#1}\right]}
\newcommand{\prob}[1]{\mathbb{P}\left[{#1}\right]}
\newcommand{\given}{\; \big\vert \;} 
\newcommand{\set}[1]{\left\{#1\right\}} 
\newcommand{\indicator}[1]{1_{\set{#1}}} 
\newcommand{\Ind}[1]{\mathbbm{1}_{#1}} 
\newcommand{\SetIn}[1]{\mathbbm{1}_{\set{#1}}} 
\newcommand{\red}[1]{\textcolor{red}{#1}} 
\newcommand{\floor}[1]{\left\lfloor#1\right\rfloor}
\newcommand{\ceil}[1]{\left\lceil#1\right\rceil}


\newcommand{\C}{\mathbb{C}}
\newcommand{\D}{\mathbb{D}}
\newcommand{\E}{\mathbb{E}}
\newcommand{\F}{\mathbb{F}}
\newcommand{\N}{\mathbb{N}}
\renewcommand{\P}{\mathbb{P}}
\newcommand{\Q}{\mathbb{Q}}
\newcommand{\R}{\mathbb{R}}
\newcommand{\Z}{\mathbb{Z}}

\newcommand{\bA}{\mathbf{A}}
\newcommand{\bI}{\mathbf{I}}
\newcommand{\bU}{\mathbf{1}}
\newcommand{\bx}{\mathbf{x}}
\newcommand{\bX}{\mathbf{X}}
\newcommand{\bZ}{\mathbf{Z}}


\newcommand{\cA}{\mathcal{A}}
\newcommand{\cB}{\mathcal{B}}
\newcommand{\cC}{\mathcal{C}}
\newcommand{\cF}{\mathcal{F}}
\newcommand{\cG}{\mathcal{G}}
\newcommand{\cM}{\mathcal{M}}
\newcommand{\cN}{\mathcal{N}}
\newcommand{\cP}{\mathcal{P}}
\newcommand{\cT}{\mathcal{T}}
\newcommand{\cX}{\mathcal{X}}
\newcommand{\cY}{\mathcal{Y}}

\newcommand{\sA}{\mathscr{A}}
\newcommand{\sB}{\mathscr{B}}
\newcommand{\sC}{\mathscr{C}}
\newcommand{\sE}{\mathscr{E}}
\newcommand{\sF}{\mathscr{F}}
\newcommand{\sG}{\mathscr{G}}
\newcommand{\sH}{\mathscr{H}}
\newcommand{\sL}{\mathscr{L}}
\newcommand{\sS}{\mathscr{S}}
\newcommand{\sT}{\mathscr{T}}
\newcommand{\sX}{\mathscr{X}}
\newcommand{\sY}{\mathscr{Y}}
\newcommand{\sZ}{\mathscr{Z}}

\newcommand{\tS}{\tilde{S}}
% Debug
\newcommand{\todo}[1]{\begin{color}{blue}{{\bf~[TODO:~#1]}}\end{color}}

% a few handy macros

\renewcommand{\le}{\leqslant}
\renewcommand{\ge}{\geqslant}
\newcommand\matlab{{\sc matlab}}
\newcommand{\goto}{\rightarrow}
\newcommand{\bigo}{{\mathcal O}}
%\newcommand{\half}{\frac{1}{2}}
%\newcommand\implies{\quad\Longrightarrow\quad}
\newcommand\reals{{{\rm l} \kern -.15em {\rm R} }}
\newcommand\complex{{\raisebox{.043ex}{\rule{0.07em}{1.56ex}} \hskip -.35em {\rm C}}}


% macros for matrices/vectors:

% matrix environment for vectors or matrices where elements are centered
\newenvironment{mat}{\left[\begin{array}{ccccccccccccccc}}{\end{array}\right]}
\newcommand\bcm{\begin{mat}}
\newcommand\ecm{\end{mat}}

% matrix environment for vectors or matrices where elements are right justifvied
\newenvironment{rmat}{\left[\begin{array}{rrrrrrrrrrrrr}}{\end{array}\right]}
\newcommand\brm{\begin{rmat}}
\newcommand\erm{\end{rmat}}

% for left brace and a set of choices
%\newenvironment{choices}{\left\{ \begin{array}{ll}}{\end{array}\right.}
\newcommand\when{&\text{if~}}
\newcommand\otherwise{&\text{otherwise}}
% sample usage:
%\delta_{ij} = \begin{choices} 1 \when i=j, \\ 0 \otherwise \end{choices}


% for labeling and referencing equations:
\newcommand{\eql}{\begin{equation}\label}
\newcommand{\eqn}[1]{(\ref{#1})}
% can then do
%\eql{eqnlabel}
%...
%\end{equation}
% and refer to it as equation \eqn{eqnlabel}.
% some useful macros for finite difference methods:
\newcommand\unp{U^{n+1}}
\newcommand\unm{U^{n-1}}
% for chemical reactions:
\newcommand{\react}[1]{\stackrel{K_{#1}}{\rightarrow}}
\newcommand{\reactb}[2]{\stackrel{K_{#1}}{~\stackrel{\rightleftharpoons}
 {\scriptstyle K_{#2}}}~}
\makeatletter
\def\th@plain{%
\thm@notefont{}% same as heading font
\itshape % body font
}
\def\th@definition{%
\thm@notefont{}% same as heading font
\normalfont % body font
}
\makeatother
\date{}



\title{ Lecture-23: Invariant Distribution}
\author{}

\begin{document}
\maketitle

\section{Invariant Distribution} 
Let $X: \Omega \to \sX^{\Z_+}$ be a time-homogeneous Markov chain with transition probability matrix $P: \sX\times\sX\to[0,1]$. 
%\begin{defn} 
%For any countable set $\sX$, we define the set of probability mass functions over $\sX$ as 
%\EQ{
%\cM(\sX) \triangleq \set{\nu \in  [0,1]^\sX: \sum_{x\in\sX}\nu_x = 1}.
%}
%\end{defn}
\begin{defn} 
A probability distribution $\pi \in \cM(\sX)$ is said to be %\textbf{stationary distribution} or 
\textbf{invariant distribution} for the Markov chain $X$ if it satisfies the \textbf{global balance equation} $\pi = \pi P$. 
%that is $\pi_y = \sum_{x \in \sX}\pi_xp_{xy}$ for all $y \in \sX$. 
\end{defn}
\begin{defn} 
When the initial distribution of a Markov chain is $\nu \in \cM(\sX)$, 
then the conditional probability is denoted by $P_\nu: \sF \to [0,1]$ defined by 
\EQ{
P_\nu(A) \triangleq \sum_{x\in\sX}\nu(x)P_x(A)\text{ for all events }A \in \sF.
} 
\end{defn}
\begin{defn} 
For a Markov chain $X:\Omega\to\sX^{\Z_+}$, 
we denote the distribution of random variable $X_n:\Omega\to\sX$ by $\nu_n \in \cM(\sX)$ for all $n \in \Z_+$. 
That is, $\nu_n(x) \triangleq P_{\nu_0}\set{X_n = x}$ for all $x \in \sX$. 
\end{defn}
\begin{rem} 
We observe that $\nu_n(x) = \sum_{z\in\sX}\nu_0(z)(P^n)_{zx}$ for all $x \in \sX$. 
\end{rem}
\begin{shaded*}
\begin{rem} 
Facts about the invariant distribution $\pi$. 
\begin{enumerate}[i\_]
\item The global balance equation $\pi = \pi P$ is a matrix equation, that is we have a collection of $\abs{\sX}$ equations $\pi_y = \sum_{x \in \sX}\pi_xp_{xy}$ for each $y \in \sX$. 
%\item Balance equation across cuts is $\pi_y(1-p_{yy}) = \pi_y \sum_{x \neq y}p_{yx} = \sum_{x \neq y}\pi_xp_{xy}$. 
%In general, for any set $A \subseteq \sX$, we have $\sum_{y\notin A}\sum_{x\in A}\pi_xp_{xy} = \sum_{x \in A}\sum_{y\notin A}\pi_yp_{yx}$. 
%To this see, we observe that by exchanging sums and exchanging $x$ and $y$ as running variables, we can write the LHS as 
%\EQ{
%\sum_{y\in A}\pi_y\sum_{x\notin A}p_{yx} 
%= \sum_{y\in A}\pi_y(1- \sum_{x\in A}p_{yx}) 
%= \sum_{y\in A}\pi_y - \sum_{x\in A}\sum_{y\in A}\pi_yp_{yx}
%= \sum_{y\in A}\pi_y - \sum_{x\in A}(\pi_x- \sum_{y\notin A}\pi_yp_{yx}).
%}
%Thus, we observe that LHS is equal to the RHS.  
\item The invariant distribution $\pi$ is left eigenvector of stochastic matrix $P$ with the largest eigenvalue $1$. 
The all ones vector is the right eigenvector of this stochastic matrix $P$ for the eigenvalue $1$.  
\item From the Chapman-Kolmogorov equation for initial probability vector $\pi$, we have $\pi = \pi P^n$ for $n \in \N$. 
That is, if $\nu_0 = \pi$, then $\nu_n = \pi$ for all $n \in \Z_+$. 
%That is, if $P\set{X_0 = x} = \pi_x$ for each $x \in \sX$, then $P_{\pi}\set{X_n = y} = \pi_y$ for each $y \in \sX$ and all $n \in \Z_+$. 
%, since  $P_{\nu}\set{X_n = y} %= \sum_{x \in \sX}P_{\nu}\set{X_0 = x}p_{xy}^{(n)} =  \sum_{x \in \sX}\nu(x)p_{xy}^{(n)}$. 
\item Resulting process with initial distribution $\pi$ is stationary, and hence have shift-invariant finite dimensional distributions. 
For example, for any $k, n \in \Z_+$ and $x_0, \dots, x_n \in \sX$, we have 
\EQ{
P_{\pi}\set{X_0 = x_0, \dots, X_n = x_n} = P_{\pi}\set{X_k = x_0, \dots, X_{k+n} = x_n} = \pi_{x_0}p_{x_0x_1} \dots p_{x_{n-1}x_n}. 
}
\item For an irreducible Markov chain, if $\pi_x > 0$ for some $x \in \sX$, then the entire invariant vector $\pi$ is positive. 
To this end, we will show that $\pi_y > 0$ for all states $y \in \sX$. 
Let $y \in \sX$, then from the irreducibility of Markov chain, there exists an $m \in \Z_+$ such that $p_{xy}^{(m)} > 0$. 
Further, $\pi = \pi P^m$ and hence $\pi_y \ge \pi_x p_{xy}^{(m)} > 0$.  

\item Any scaled version of $\pi$ satisfies the global balance equation. 
Therefore, for any invariant vector $\alpha \in \sX^{\R_+}$ of a positive recurrent transition matrix $P$,  
the sum $\norm{\alpha}_1 = \sum_{x \in \sX}\alpha_x$ must be finite. 
We can normalize $\alpha$ and get an invariant probability measure $\pi = \frac{\alpha}{\norm{\alpha}_1}$. 
\end{enumerate}
\end{rem}
\end{shaded*}

\begin{thm} 
An irreducible Markov chain with transition probability matrix $P:\sX\times\sX\to[0,1]$ is positive recurrent iff there exists a unique invariant probability measure $\pi \in \cM(\sX)$ that satisfies global balance equation $\pi = \pi P$ and $\pi_x = \frac{1}{\mu_{xx}} > 0$ for all $x \in \sX$. 
\end{thm}
\begin{proof} 
Consider an irreducible Markov chain $X:\Omega\to\sX^{\Z_+}$ with transition probability matrix $P$.  
We will first show that positive recurrence of $X$ implies the existence of a positive invariant distribution $\pi$ and its uniqueness. 
Then we will show that the existence of a unique positive invariant distribution $\pi$ implies positive recurrence of $X$.  
\begin{itemize}
\item[Implication:]
For Markov chain $X$, let the initial state be $X_0 = x$. 
Recall that the number of visits to state $y \in \sX$ in the first $n$ steps of the Markov chain $X$ is denoted by $N_y(n) = \sum_{k=1}^n\SetIn{X_k = y}$.  
It follows that $\sum_{y \in \sX}N_y(n) = n$ for each $n \in \N$. 
Let $H_x \triangleq \tau_X^{\set{x},1}$ be the first recurrence time to state $x\in\sX$, then we have $N_x(H_x) = 1$ and $\sum_{y \in \sX}N_y(H_x) = H_x$. 

\begin{itemize}
\item[Existence:] 
We define a vector $v \in \R^\sX$ by $v_y \triangleq \E_x[N_y(H_x)]$ for each $y \in \sX$. 
We observe that $v_y \ge 0$ for each state $y \in \sX$. 
In particular, $v_x = 1$. 
Since $X$ is positive recurrent, we get that $\sum_{y \in \sX}v_y = \E_xH_x = \mu_{xx} < \infty$. 
We will show that the vector $v$ satisfies the global balance equation $v = vP$. 
To see this, we first define $\lambda_{xy}^{(n)}  \triangleq P_x\set{X_n = y, n \le H_x}$ for all $n \in \N$ and states $x,y\in\sX$. 
We observe that $\lambda_{xy}^{(1)} = p_{xy}$ for each $y \in \sX$. 
Next, we observe from the monotone convergence theorem, that 
\EQN{
\label{eqn:Eigenvector}
v_y = \E_x N_y(H_x) = \E_x\sum_{n \in \N}\SetIn{X_n = y, n \le H_x} = \sum_{n \in \N}P_x\set{X_n = y, n \le H_x} = \sum_{n\in\N} \lambda_{xy}^{(n)}. 
}
For $n \ge 2$, partitioning the event $\set{X_n=y, n \le H_x}$ by the events $(\set{X_{n-1}=z}: z \in \sX\setminus\set{x})$, the countable additivity of conditional probability for disjoint events, and the definition of conditional probability, we get 
\EQ{
\lambda_{xy}^{(n)} 
%= \sum_{z \neq x}P_x\set{X_n = y, X_{n-1} = z, n \le H_x}\\
= \sum_{z \neq x}P(\set{X_n = y}\given\set{X_{n-1} = z, n \le H_x, X_0 = x}) P_x\set{X_{n-1} = z, n \le H_x}.
}
Recall that since $H_x$ is adapted to natural filtration $\sF_\bullet$ of Markov chain $X$, 
we have $\set{n \le H_x, X_0=x} = \set{X_0=x}\cap\set{H_x > n-1}^c \in \sF_{n-1}$. 
Together with the Markov property of $X$ and the fact that $\set{X_{n-1}=z, n \le H_x} = \set{X_{n-1}=z, n-1 \le H_x}$, we obtain 
\EQN{
\label{eqn:EigenvectorTerms}
\lambda_{xy}^{(n)} =  \sum_{z \neq x}P(\set{X_n = y}\given\set{X_{n-1} = z}) P_x\set{X_{n-1} = z, n-1 \le H_x} = \sum_{z \neq x}\lambda_{xz}^{(n-1)}p_{zy}.
} 
%From the definition of $\lambda_{xy}^{(n)}$, we have $v_y = \sum_{n \in \N}\lambda_{xy}^{(n)}$ for each $y \in \sX$.  
Substituting the expression for $\lambda_{xy}^{(n)}$ in~\eqref{eqn:EigenvectorTerms} into the expression for  $v_y = \sum_{n \in \N}\lambda_{xy}^{(n)}$ in~\eqref{eqn:Eigenvector} and using the fact that $v_x = 1$, we obtain 
\EQ{
v_y = p_{xy} +  \sum_{n \ge 2}\sum_{z \neq x}\lambda_{xz}^{(n-1)}p_{z y} 
%= p_{xy} + \sum_{z \neq x}p_{z y}\sum_{n \in \N}P_x\set{X_n = z, n \le H_x} 
= v_xp_{xy} +  \sum_{z \neq x}v_{z}p_{z y} = \sum_{x \in \sX}v_xp_{xy}. 
}
Since $v$ has a finite sum, it follows that $\pi \triangleq \frac{v}{\sum_{x \in \sX}v_x}$ is an invariant distribution for the Markov chain $X$ with the transition matrix $P$. 
In addition, we have $\pi_x = \frac{v_x}{\sum_{y \in \sX}v_y} = \frac{1}{\mu_{xx}} > 0$. 

\item[Uniqueness:] Next, we show that this is a unique invariant measure independent of the initial state $x$, 
and hence $\pi_y = \frac{1}{\mu_{yy}} > 0$ for all $y \in \sX$.  
For uniqueness, we observe from the Chapman-Kolmogorov equations and invariance of $\pi$ that $\pi = \frac{1}{n}\pi(P + P^2 + \dots + P^n)$.  
Hence, $\pi_y = \sum_{x \in \sX}\pi_x \frac{1}{n}\sum_{k=1}^np_{xy}^{(k)}$ for all states $y \in \sX$.
%For uniqueness, we observe that for any initial distribution $\nu$
%\EQ{
%\lim_{n \to \infty}P_\nu\set{X_n = y}=\lim_{n \to \infty}\sum_{x \in \sX}\nu(x)\frac{1}{n}\sum_{k=1}^np_{xy}^{(k)}.
%}
Taking limit $n \to \infty$ on both sides, and exchanging limit and summation on right hand side using bounded convergence theorem for summable series $\pi$, we get $\pi_y = \frac{1}{\mu_{yy}}\sum_{x \in \sX}\pi_x =  \frac{1}{\mu_{yy}} > 0$ for all states $y \in \sX$.  
\end{itemize}
\item[Converse:]
Let $\pi$ be the unique positive invariant distribution of Markov chain $X$, such that $\pi_y = \frac{1}{\mu_{yy}} > 0$ for all states $y \in \sX$. 
It follows that the Markov chain $X$ is positive recurrent. 
\end{itemize}
\end{proof}

\begin{cor} 
An irreducible Markov chain on a finite state space $\sX$ has a unique positive stationary distribution $\pi$. 
\end{cor}
%\begin{proof} 
%Let $X$ be a positive recurrent Markov chain on state space $\sX$. 
%Let $\nu_x(n)$ be the probability of the process $X$ being in a state $x$ at time $n \in \Z_+$. 
%From bounded convergence theorem, we have 
%\EQ{
%\lim_{n \to \infty}\frac{1}{n}\sum_{k=1}^n\nu_y(k) = \lim_{n \to \infty}\sum_{x \in \sX}\nu_x(0) \frac{1}{n}\sum_{k=1}^{n}p_{xy}^{(k)} =  \sum_{x \in \sX}\nu_x(0) \lim_{n \to \infty}\frac{1}{n}\sum_{k=1}^{n}p_{xy}^{(k)} = \sum_{x \in \sX}\nu_x(0)\frac{1}{\mu_{yy}} = \frac{1}{\mu_{yy}} > 0. 
%}
%This implies $\lim_{n \to \infty}\nu(n) = (\frac{1}{\mu_{yy}}: y \in \sX)$. 
%Let $\nu$ be the vector such that $\nu_x = \frac{1}{\mu_{xx}}$. 
%We also observe that $\lim_{n \to \infty}\frac{1}{n}\sum_{k=1}^n\nu_y(k-1) = \frac{1}{\mu_{yy}} = \nu_y$. 
%Since $\frac{1}{n}\sum_{k=1}^n\nu_y(k) = \sum_{x \in \sX}p_{xy}\frac{1}{n}\sum_{k=1}^n\nu_x(k-1)$, 
%taking limits on both sides using bounded convergence theorem, 
%we obtain the global balance equation $\nu = \nu P$. 
%We need to show that $\nu$ is summable and unique. 
%For summability, we have for any finite set $F \subseteq S$, 
%\EQ{
%\sum_{i \in F}\nu_x = \sum_{i \in F}\lim_{n \to \infty}\nu_x(n) = \lim_{n \to \infty}\sum_{i \in F}\nu_x(n) \le 1. 
%} 
%This implies that $0 \le \sum_{x \in \sX}\nu_x \le 1$,  and we let the invariant distribution be $\pi = \frac{1}{\sum_{x \in \sX}\nu_x}\nu$. 
%For uniqueness, let $\pi$ be the invariant distribution of probability transition matrix $P$. 
%Hence,
%\EQ{
%\pi = \frac{1}{n}\sum_{k=1}^{n}\pi = \pi\frac{1}{n}\sum_{k=1}^nP^k. 
%}
%Taking limits on both sides, we get $\pi = \nu$. 
%
%Conversely, let $\pi$ be the invariant probability vector. 
%Then, if the Markov chain was transient or null recurrent, we would have $\lim_{n \to \infty}\frac{1}{n}\sum_{k=1}^np_{xy}^{(k)} = 0$. 
%Since $\pi$ is an invariant vector, we get $\pi = \pi P^k$ for each $k \in \N$ and hence $\pi = \pi\frac{1}{n}\sum_{k=1}^nP^k$. 
%Taking limit on both sides, we have $\pi > 0$ being equal to zero on right hand side, yielding a contradiction. 
%\end{proof}
\begin{defn}
An irreducible, aperiodic, positive recurrent Markov chain is called \textbf{ergodic}. 
\end{defn}
\begin{shaded*}
\begin{rem} 
Additional remarks about the stationary distribution $\pi$. 
\begin{enumerate}[i\_]
\item For a Markov chain with multiple positive recurrent communicating classes $\cC_1, \dots, \cC_m$, 
one can find the positive equilibrium distribution for each class, 
and extend it to the entire state space $\sX$ denoting it by $\pi_k$ for class $k \in [m]$.  
It is easy to check that any convex combination $\pi = \sum_{k = 1}^m\alpha_k\pi_k$ satisfies the global balance equation $\pi = \pi P$, 
where $\alpha_k \ge 0$ for each $k \in [m]$ and $\sum_{k=1}^m\alpha_k = 1$. 
Hence, a Markov chain with multiple positive recurrent classes have a convex set of invariant probability measures, 
with the individual invariant distribution $\pi_k$ for each positive recurrent class $k \in [m]$ being the extreme points. 

\item Let $\mu(0) = e_x$, that is let the initial state of the positive recurrent Markov chain be $X_0 = x$. 
Then, we know that %for an aperiodic Markov chain $X$, 
\EQ{
\pi_y = \frac{1}{\mu_{yy}} = \lim_{n \to \infty}\frac{1}{n}\sum_{k=1}^np_{xy}^{(k)} =  \lim_{n \to \infty}\frac{1}{n}\E_xN_y(n).
}
That is, $\pi_y$ is limiting average of number of visits to state $y \in \sX$. 
\item If a positive recurrent Markov chain is aperiodic, then limiting probability of being in a state $y$ is its invariant probability, 
that is $\pi_y = \lim_{n \to \infty}p_{xy}^{(n)}$. 
\end{enumerate}
\end{rem}
\end{shaded*}
\red{
\begin{thm} 
For an ergodic Markov chain $X$ with invariant distribution $\pi$, and $n$th step distribution $\mu(n)$, we have $\lim_{n \to \infty}\mu(n) = \pi$ in the total variation distance. 
\end{thm}
\begin{proof} 
Consider independent time homogeneous Markov chains $X: \Omega \to \sX^{\Z_+}$ and $Y: \Omega \to \sX^{\Z_+}$ each with transition matrix $P$. 
The initial state of Markov chain $X$ is assumed to be $X_0 = x$, whereas the Markov chain $Y$ is assumed to have an initial distribution $\pi$. 
It follows that $Y$ is a stationary process, while $X$ is not. 
In particular, 
\meq{2}{
&\mu_y(n) = P_x\set{X_n = y} = p_{xy}^{(n)}, && P_{\pi}\set{Y_n = y} = \pi_y.
} 
Let $\tau = \inf\{n \in \Z_+: X_n = Y_n\}$ be the first time that two Markov chains meet, called the \textbf{coupling time}. 
\begin{itemize}
\item[Finiteness:] First, we show that the coupling time is almost surely finite. 
To this end, we define a a new Markov chain on state space $\sX \times \sX$ with transition probability matrix $Q$ such that $q((x,w), (y,z)) = p_{xy}p_{wz}$ for each pair of states $(x,w), (y,z) \in \sX \times \sX$. 
The $n$-step transition probabilities for this couples Markov chain are given by 
\EQ{
q^{(n)}((x,w), (y,z)) \triangleq p^{(n)}_{xy}p^{(n)}_{wz}.
} 
\begin{itemize}
\item[Ergodicity:] 
Since the Markov chain $X$ with transition probability matrix $P$ is irreducible and aperiodic, 
for each $x,y,w,z \in \sX$ there exists an $n_0 \in \Z_+$ such that $q^{(n)}((x,w), (y,z)) = p^{(n)}_{xy}p^{(n)}_{wz} > 0$ for all $n \ge n_0$ from a previous Lemma on aperiodicity.   
Hence, the irreducibility and aperiodicity of this new \textbf{product} Markov chain follows. 
\item[Invariant:]
It is easy to check that $\theta(x,w) = \pi_x\pi_{w}$ is the invariant distribution for this product Markov chain, 
since $\theta(x,w) > 0$ for each $(x,w) \in \sX \times \sX$, $\sum_{x,w \in \sX}\theta(x,w) = 1$, and for each $(y,z) \in \sX \times \sX$, we have 
\EQ{
\sum_{x,w \in \sX}\theta(x,w)q((x,w), (y,z)) = \sum_{x \in \sX}\pi_xp_{xy}\sum_{w \in \sX}\pi_{w}p_{wz} = \pi_y\pi_{z} = \theta(y,z). 
}
\item[Recurrence:] 
This implies that the product Markov chain is positive recurrent, and each state $(x,x) \in \sX \times \sX$ is reachable with unit probability from any initial state $(y,w) \in \sX \times \sX$. 
\end{itemize}
In particular, the coupling time is almost surely finite.  
\item[Coupled process:] Second, we show that from the coupling time onwards, the evolution of two Markov chains is identical in distribution.  
That is, for each $y \in \sX$ and $n \in \Z_+$, 
\EQ{
P_{X_{\tau}}\set{X_{n} = y, n \ge \tau} = P_{Y_{\tau}}\set{Y_{n} = y, n \ge \tau}. 
}
This follows from the strong Markov property for the joint process where $\tau$ is stopping time for the joint process $((X_n, Y_n): n \in \Z_+)$ such that $X_\tau = Y_\tau$, 
and both marginals  have the identical transition matrix. % and that $X_{\tau} = Y_{\tau}$. 
\item[Limit:] For any $y \in \sX$, we can write the difference as 
\EQ{
\abs{p_{xy}^{(n)} - \pi_y} = \abs{P_x\set{X_n = y, n < \tau} - P_{\pi}\set{Y_n = y, n < \tau}} \le 2P_{\delta_x, \pi}(\tau > n). 
}
Since the coupling time is almost surely finite for each initial state $x,y \in \sX$, 
we have $\sum_{n \in \N}P_{\delta_x,\pi}\set{\tau = n} = 1$ and the tail-sum $P_{\delta_x,\pi}\set{\tau > n}$ goes to zero as $n$ grows large, and the result follows. 
\end{itemize}
\end{proof}
}
\section{Computing invariant distribution}
When the state space $\sX$ is finite, one can find left eigenvector of probability transition matrix $P$ for the largest eigenvalue~$1$. 
This is the invariant distribution that satisfies the global balance equation $\pi = \pi P$. 
\begin{defn} 
Consider a time homogeneous positive recurrent Markov chain $X:\Omega\to\sX^{\Z_+}$ with probability transition matrix $P$ and invariant distribution $\pi \in \cM(\sX)$. 
For any two disjoint sets $A, B \subseteq \sX$,  
the probability flux from set of nodes $A$ to set of nodes $B$ is defined as $\Phi(A,B) =\sum_{x \in A}\sum_{y \in B}\pi_xp_{xy}$. 
\end{defn}
\begin{rem} 
The probability flux from a single node $x$ to single node $y$ is denoted by $\Phi(x,y) = \pi_xp_{xy}$. 
\end{rem}
\begin{defn} 
For a time homogeneous Markov chain $X:\Omega\to\sX^{\Z_+}$ with probability transition matrix $P$ 
represented as the weighted transition graph $G = (\sX, E, w)$, 
a \textbf{cut} is defined as the partition $(\sX_1,\sX_2)$ of nodes.
\end{defn}
\begin{lem}
Probability flux balances across cuts. 
\end{lem}
\begin{proof}
A cut for the state space $\sX$ is given by a partition $(\sX_1,\sX_2)$. 
We show that $\Phi(\sX_1,\sX_2) = \Phi(\sX_2,\sX_1)$. 
To this see, we observe that by exchanging sums from the monotone convergence theorem, and exchanging $x$ and $y$ as running variables, we can write the probability flux $\Phi(\sX_1,\sX_2)$ as 
\EQ{
\sum_{y\in \sX_1}\pi_y\sum_{x\notin\sX_1}p_{yx} 
= \sum_{y\in\sX_1}\pi_y(1- \sum_{x\in\sX_1}p_{yx}) 
= \sum_{y\in\sX_1}\pi_y - \sum_{x\in\sX_1}\sum_{y\in\sX_1}\pi_yp_{yx}
= \sum_{y\in\sX_1}\pi_y - \sum_{x\in\sX_1}(\pi_x- \sum_{y\notin\sX_1}\pi_yp_{yx}).
}
%Then, we observe that 
%\eq{
%\Phi(\sX_1,\sX_2) 
%&= \sum_{x \in \sX_1}\sum_{y \in \sX_2}\pi_xp_{xy}
%= \sum_{y \in \sX_2}\sum_{x \in \sX}\pi_xp_{xy} - \sum_{y \in \sX_2}\sum_{x\in \sX_2}\pi_xp_{xy}\\
%&= \sum_{x \in \sX_2}\pi_x\sum_{y \in \sX}p_{xy} - \sum_{x \in \sX_2}\sum_{y\in \sX_2}\pi_xp_{xy}
%= \sum_{x \in \sX_2}\sum_{y \in \sX_1}\pi_xp_{xy}
%= \Phi(\sX_2, \sX_1).
%}
\end{proof}
\begin{cor}
For any states $y \in \sX$, we have $\pi_y(1-p_{yy}) = \pi_y \sum_{x \neq y}p_{yx} = \sum_{x \neq y}\pi_xp_{xy}$. 
\end{cor}
\begin{proof}
It follows from probability flux balancing across the cut $(\set{y}, \set{y}^c)$. 
\end{proof}
%\begin{lem}
%Consider the state space $\sX = \Z_+$ a function $f:\sX\to[0,1]$, and the transition probability $p_{xy} = f(\abs{x-y})\SetIn{x \neq y}+(1-\sum_{y \neq x}f(\abs{x-y}))\SetIn{x=y}$ for all states $x,y\in\Z$, 
%then 
%\EQ{
%\Pi(z) = \frac{\sum_{x\in \Z_+}z^x\pi_x\Big(\sum_{w=0}^{x}z^{-w}f(w) -\sum_{w=0}^xf(w)\Big)}{\sum_{w\in\Z_+}f(w)-F(z)}.
%}
%\end{lem}
%\begin{proof}
%Defining probability generating function for the invariant distribution and function $f$ as $\Pi(z) \triangleq \sum_{y\in\Z_+}z^y\pi_y$ and $F(z) \triangleq \sum_{w\in\Z_+}z^wf(w)$ respectively, we get 
%\EQ{
%\Pi(z)
%= \sum_{x\in \Z_+}z^x\pi_x\sum_{y\in \Z_+}z^{y-x}p_{xy}
%= \sum_{x\in \Z_+}z^x\pi_x\Big(\sum_{w\in\Z_+}z^{w}f(w) +\sum_{w=0}^{x}z^{-w}f(w) + 1-\sum_{w\in\Z_+}f(w)-\sum_{w=0}^xf(w)\Big).
%}
%\end{proof}

%\begin{shaded*}
%\begin{exmp}[Single Server Queue]
%\end{exmp}
%\end{shaded*}

\end{document}
